% ED26-SRC-063-0104: original pp.247--307; this packet treats §§1--3, pp.247--265.
\chapter{Aristotle's Modal Syllogisms: Johnson's Competing Reconstructions}
\label{sec:p13-src063-0104}
\chaptermark{Aristotle's Modal Syllogisms}

Fred Johnson's 2004 chapter asks how formal systems can match Aristotle's
judgments of validity and invalidity, including modal syllogisms.
An assertion calculus without rejection and countercases would answer only
half that question. His sequence begins with \L{}ukasiewicz's assertoric
system \L{}A, examines the four-valued modal system \L{}M and modern
modal-predicate translations, and then turns to McCall's L-X-M and
contingency calculi. These are materially different reconstructions, not
variants interchangeable under the title Aristotle
\parencite[pp.~247--265]{Johnson2004ModalSyllogisms}.

\section{Historical interpretation and the assertoric language}
\label{sec:p13-johnson-assertoric}

\L{}ukasiewicz reads Barbara as a conditional claim, rather than an
inferential discourse with two premises and a conclusion. Johnson notes
that Smith's translation of \emph{Prior Analytics} 25b37--40 can seem to
support this reading, while Corcoran and Smiley defend natural-deduction
interpretations. An equivalent modern conditional does not by itself
settle what Aristotle took a syllogism to be
\parencite[pp.~247--248]{Johnson2004ModalSyllogisms}.

Johnson replaces Polish notation with contemporary connective notation.
That change is explicitly his presentation choice.
Term variables are \(a,b,c,\ldots\), atomic categorical formulas are
\(Aab,Iab\), and formulas are closed under \(\neg\) and \(\to\).
There are no other formulas except those generated in this way.
Here \(Aab\) means all \(a\) are \(b\), and \(Iab\) some \(a\) are \(b\):
the subject is first. Define
\[
 Eab\mathrel{:\equiv}\neg Iab,\quad
 Oab\mathrel{:\equiv}\neg Aab,\quad
 (p\land q)\mathrel{:\equiv}\neg(p\to\neg q),\quad
 (p\leftrightarrow q)\mathrel{:\equiv}(p\to q)\land(q\to p).
\]
Boger's \(AaB\), with its predicate first, must therefore not be
identified with Johnson's \(Aab\) by keeping letter positions fixed.

Johnson uses \(\vdash p\) for assertion and \(\dashv p\) for rejection.
They are relative to the indicated system. Rejection of \(p\) is not
assertion of \(\neg p\): \(Aab\) is not generally a thesis, but neither is
its negation. Failure of universal validity is distinct from universal
falsity \parencite[pp.~248--249]{Johnson2004ModalSyllogisms}.

The assertion basis consists of every classical propositional tautology
(A0), \(Aaa\) (A1), \(Iaa\) (A2), Barbara
\(Abc\to(Aab\to Aac)\) (A3), and Datisi
\(Abc\to(Iba\to Iac)\) (A4). Assertion transformations are uniform
term substitution, modus ponens, and interchange of a definiendum with
its definiens. Substitution may identify different variables but cannot
be inverted: an assertion for a special identification does not establish
an assertion for arbitrary distinct terms.

A2 carries existential force. A modern extensional comparison takes every
term extension nonempty; otherwise \(Iaa\) need not hold. It is not enough
to say only that the universe itself is nonempty. These assumptions also
matter for subalternation and Darapti. They are not silently imposed on
every other historical system.

\begin{hbtheorem}{thm:p13-johnson-la-immediate}{proved=annex:p13.johnson.immediate}{Immediate assertoric inferences in Johnson's assertion basis}
In the assertion basis just specified, \(Iab\to Iba\), \(Eab\to Eba\),
\(Aab\to Iab\), \(Eab\to Oab\), \(Aab\to Iba\), and \(Eab\to Oba\)
are asserted. The oppositions
\(\neg Aab\leftrightarrow Oab\),
\(\neg Eab\leftrightarrow Iab\),
\(\neg Iab\leftrightarrow Eab\), and
\(\neg Oab\leftrightarrow Aab\) are also asserted.
\end{hbtheorem}

Johnson's derived assertion operations are reversal (the contraposition
forms of a conditional and of a two-antecedent conditional), antecedent
interchange, antecedent strengthening, consequent weakening, biconditional
formation from both implications, and substitution of equivalents.
Each relies on propositional tautologies plus modus ponens; replacement
inside arbitrary formulas additionally uses induction on formula
formation. They are assertion operations, not instructions for reversing
every arbitrary object-language argument
\parencite[pp.~250--252]{Johnson2004ModalSyllogisms}.

\section{The complete assertoric mood comparison}
\label{sec:p13-johnson-moods}

The following is Johnson's table~1 in subject-first notation, with the
assertion mark understood in every row. Each pair in the last column
denotes antecedents followed by the consequent of the nested conditional.
It is a table of source claims, not a new editorial rule display.
\parencite[p.~253]{Johnson2004ModalSyllogisms}

\begin{longtable}{@{}llp{0.54\textwidth}@{}}
\caption{The complete assertoric mood comparison}\\
\toprule
Figure & Name & Antecedents; consequent\\
\midrule
\endfirsthead
\caption[]{The complete assertoric mood comparison (continued)}\\
\toprule Figure & Name & Antecedents; consequent\\ \midrule
\endhead
\midrule\multicolumn{3}{r}{\emph{continued on next page}}\\
\endfoot
\endlastfoot
1 & Barbara & \(Abc,Aab;Aac\)\\
1 & Celarent & \(Ebc,Aab;Eac\)\\
1 & Darii & \(Abc,Iab;Iac\)\\
1 & Ferio & \(Ebc,Iab;Oac\)\\
2 & Cesare & \(Ecb,Aab;Eac\)\\
2 & Camestres & \(Acb,Eab;Eac\)\\
2 & Festino & \(Ecb,Iab;Oac\)\\
2 & Baroco & \(Acb,Oab;Oac\)\\
3 & Darapti & \(Abc,Aba;Iac\)\\
3 & Felapton & \(Ebc,Aba;Oac\)\\
3 & Disamis & \(Ibc,Aba;Iac\)\\
3 & Datisi & \(Abc,Iba;Iac\)\\
3 & Bocardo & \(Obc,Aba;Oac\)\\
3 & Ferison & \(Ebc,Iba;Oac\)\\
4 & Bramantip & \(Acb,Aba;Iac\)\\
4 & Camenes & \(Acb,Eba;Eac\)\\
4 & Dimaris & \(Icb,Aba;Iac\)\\
4 & Fresison & \(Ecb,Iba;Oac\)\\
4 & Fesapo & \(Ecb,Aba;Oac\)\\
Subaltern & Barbari & \(Abc,Aab;Iac\)\\
Subaltern & Celaront & \(Ebc,Aab;Oac\)\\
Subaltern & Cesaro & \(Ecb,Aab;Oac\)\\
Subaltern & Camestrop & \(Acb,Eab;Oac\)\\
Subaltern & Camenop & \(Acb,Eba;Oac\)\\
\bottomrule
\end{longtable}

Barbara and Datisi are the two categorical axiom patterns. The other
assertions are generated by the immediate inferences and propositional
operations. For example, Datisi followed by antecedent interchange and
particular conversion yields Disamis. Reversal of Disamis, rewriting the
negative particulars as \(E\), and a uniform exchange of term letters
yields Celarent. Darii replaces the converted particular antecedent of
Datisi through strengthening. The subalterns weaken universal conclusions
to particulars; this last step uses the existential basis.
The source's fourth-figure and subaltern organization is not Boger's
three-figure definition of a syllogism.

\section{Rejection and contrasted instances}
\label{sec:p13-johnson-rejection}

An elementary formula is a chain
\(x_1\to(x_2\to\cdots\to x_n)\), with every \(x_i\) a simple
categorical \(Apq,Ipq,Epq,Opq\). Johnson adds rejection axiom R1,
\[
 \dashv(Acb\to(Aab\to Iac)).
\]
Uniform substitution works backward for rejection: if a substitution
instance is rejected, the more general formula is rejected.
Rejection detachment uses an asserted \(x\to y\) and rejected \(y\)
to reject \(x\). Slupecki's additional operation combines rejected
\(x\to z\) and \(y\to z\) into rejected \(x\to(y\to z)\), but only
when \(x,y\) are negated positive categoricals and \(z\) has the required
elementary form. Definition interchange is also allowed.
The side conditions are part of the operation; dropping them changes
the calculus \parencite[pp.~252--254]{Johnson2004ModalSyllogisms}.

Rejection reversal and antecedent interchange use asserted propositional
equivalences and rejection detachment. Rejection antecedent weakening is
dual in direction to assertion strengthening: replacing a premise by a
weaker one cannot repair the rejected implication. Rejection consequent
strengthening likewise cannot repair it. Equivalent replacement and
rejection implication introduction complete the derived toolkit.
The last operation combines asserted \(p\) with rejected \(q\) to reject
\(p\to q\). These operations are about failure to be a thesis, not an
unrestricted rule that one may infer any negative assertion.

For \emph{Prior Analytics} 26a2--9, the premise pair is \(Aba,Ecb\).
Aristotle's contrasted instances use animal/man/horse and
animal/man/stone. In the first, all horses are animals while no horses
are men, so neither a universal nor a particular negative conclusion
about horses and animals is necessary. In the second, no stones are
animals while no stones are men, so neither positive conclusion is
necessary. A comparison model must make the subject kinds nonempty;
otherwise the particular-conclusion checks change
\parencite[pp.~255--257]{Johnson2004ModalSyllogisms}.

\begin{hbtheorem}{thm:p13-johnson-contrasted}{proved=annex:p13.johnson.contrasted}{Four conclusions defeated by two contrasted models}
Under nonempty-term extensional semantics, the premises \(Aba,Ecb\)
entail none of \(Ica,Oca,Aca,Eca\). Two finite models suffice: one
with nonempty \(c\subseteq a\), another with nonempty \(c\) disjoint
from \(a\), in both of which \(b\subseteq a\) and \(c\) is disjoint
from \(b\).
\end{hbtheorem}

Johnson's syntactic route first rejects AEI-1 and AEO-1 (theorems~1.6--1.7),
then rejects AEA-1 and AEE-1 by consequent strengthening (theorem~1.8).
This parallels the two-stage semantic argument: defeat the particular
conclusions, then use universal-to-particular subalternation to defeat
the stronger conclusions. Ross criticizes the absence of a traditional
distribution-rule explanation; \L{}ukasiewicz objects to logic depending
on concrete empirical vocabulary. Johnson instead treats the instances
as legitimate consistency/countermodel evidence. This is precisely a
disagreement that must remain visible beside Boger's terminology.

At 27b12--23, Johnson reads Aristotle's further technique as rejection
by premise weakening, not mere ambiguity. For \(Enm,Oxm\), the
black/snow/animal instance defeats \(Ixn,Axn\). The remaining
\(Oxn,Exn\) are defeated by observing that \(Exm\) is stronger than
\(Oxm\) and that the stronger premise pair already has countercases.
The syntax uses rejection weakening and strengthening (theorem~1.9).
A three-antecedent AAAA rejection, theorem~1.10, shows why the treatment
cannot stop at familiar two-premise candidates
\parencite[pp.~257--258]{Johnson2004ModalSyllogisms}.

Johnson explicitly illustrates, rather than fully reproduces, the general
reduction of every formula to elementary ones. His claim of
\enquote{100\% Aristotelicity} for \L{}A concerns agreement with
Aristotle's classifications of two-premise assertoric syllogisms.
It is not an accuracy percentage measured on participants, a proof of
unique historical interpretation, or automatic coverage of every modal
combination.

\section{The four-valued modal approach and its failure of fit}
\label{sec:p13-johnson-lm}

Johnson calls \L{}ukasiewicz's modal-propositional reconstruction
\L{}M. \(M\) and \(L\) mean possibility and necessity. Its truth
values are \(1,2,3,4\), with only \(1\) designated.
Rows of the implication matrix are antecedent values; columns are
consequent values. The table is part of this reconstruction, not a
truth table for Aristotle's original terms.
\parencite[pp.~258--260]{Johnson2004ModalSyllogisms}

\begin{tabular}{@{}c|cccc|ccc@{}}
 & \multicolumn{4}{c|}{\(\to\)} & \(\neg\) & \(M\) & \(L\)\\
 &1&2&3&4&&&\\ \hline
1&1&2&3&4&4&1&2\\
2&1&1&3&3&3&1&2\\
3&1&2&1&2&2&3&4\\
4&1&1&1&1&1&3&4
\end{tabular}

The unrestricted laws
\((p\to q)\to(Mp\to Mq)\) and
\((p\to q)\to(Lp\to Lq)\) are designated under this table.
Johnson finds a different reading of the ancient passages at
34a5--7, 34a22--24 and 34a29--31 more natural: the inference depends
on necessary implication, represented by
\(L(p\to q)\to(Mp\to Mq)\) and
\(L(p\to q)\to(Lp\to Lq)\). All four formulas are nevertheless
designated in this matrix. Therefore that alternative phrasing cannot
by itself repair its excessive strength.

\begin{hbtheorem}{thm:p13-johnson-lm-laws}{proved=annex:p13.johnson.lm}{Matrix implication laws and modal propagation}
Every assignment of values \(1,2,3,4\) to \(p,q\) gives designated
value \(1\) to all four implication laws just stated.
Also, for every assignment to \(p,q,r\), both
\((p\to(q\to r))\to(p\to(Lq\to Lr))\) and
\((p\to(q\to r))\to(Lp\to(q\to Lr))\) are designated.
\end{hbtheorem}

The strong \(L\)-law propagates through a two-antecedent conditional,
yielding both
\[
 (p\to(q\to r))\to(p\to(Lq\to Lr)),\qquad
 (p\to(q\to r))\to(Lp\to(q\to Lr)).
\]
Johnson's theorem~2.2 derives these by propositional manipulation of
the modal law. Consequently plain syllogisms acquire modal variants
that his reconstruction of Aristotle rejects: Camestres LXL,
Baroco LXL, Barbara XLL, and Ferio XLL. Here a sequence such as XLL
records first premise, second premise, conclusion; X is assertoric,
L necessary. Agreement on plain syllogisms is thus insufficient evidence
for an adequate modal reconstruction.

\section{Modern modal-predicate translations}
\label{sec:p13-johnson-modal-translations}

A translation must specify where a modal operator falls relative to
quantification and predication. Johnson's table~3 lists seven
Becker-type choices. All rows below pair a proposed necessary universal
with a proposed necessary particular; none is automatically the ancient
meaning \parencite[pp.~260--262]{Johnson2004ModalSyllogisms}.

\begin{tabularx}{\textwidth}{@{}l X X@{}}
\toprule & Universal & Particular\\ \midrule
(i)&\(\forall x(Ax\to\Box Bx)\)&\(\exists x(Ax\land\Box Bx)\)\\
(ii)&\(\Box\forall x(Ax\to Bx)\)&\(\Box\exists x(Ax\land Bx)\)\\
(iii)&\(\forall x\Box(Ax\to Bx)\)&\(\exists x\Box(Ax\land Bx)\)\\
(iv)&\(\forall x(\Box Ax\to\Box Bx)\)&\(\exists x(\Box Ax\land\Box Bx)\)\\
(v)&\(\forall x(\Diamond Ax\to\Box Bx)\)&\(\exists x(\Box Ax\land\Box Bx)\)\\
(vi)&\(\forall x(\Diamond Ax\to Bx)\)&\(\exists x(\Box Ax\land Bx)\)\\
(vii)&\(\forall x(\Box Ax\to Bx)\)&\(\exists x(\Box Ax\land Bx)\)\\
\bottomrule
\end{tabularx}

Johnson illustrates both directions of failure. One reading makes
Bocardo LXL valid by moving an existing witness from \(B\) to \(A\)
without changing its necessary non-\(C\) property. Another makes
Bocardo LLL fail in an S5 model, although Johnson reports Aristotle's
acceptance and ecthesis proof at 30a6--14.

\begin{hbtheorem}{thm:p13-johnson-modal-countercases}{proved=annex:p13.johnson.modal}{Three exact translation tests}
In ordinary quantified modal reasoning,
\[
 \exists x(Bx\land\Box\neg Cx),\quad
 \forall x(Bx\to Ax)
 \quad\text{entail}\quad
 \exists x(Ax\land\Box\neg Cx).
\]
In constant-domain S5 with two accessible worlds, however,
\[
 \forall x(Cx\to\Box Bx),\quad
 \exists x(\Box Ax\land\Box\neg Bx)
 \quad\text{need not entail}\quad
 \exists x(\Box Ax\land\Box\neg Cx).
\]
These are tests of the stated modern translations, not proofs of
Aristotle's own acceptance or rejection.
The mixed Darii translation
\(\forall x(Bx\to\Box Cx),\exists x(Ax\land Bx)\) likewise need
not entail \(\exists x\Box(Ax\land Cx)\) in constant-domain S5.
\end{hbtheorem}

The first proof needs ordinary witness reasoning, not a modal
distribution principle. The second fixes both the frame and domain;
changing them without argument would change the comparison.
Johnson reads ecthesis as essentially existential instantiation, while
recording alternatives from Thom and Smith. His source text has an
inconsistent Bocardo-LXL classification sentence on p.~263, contrasted
with pp.~261 and 264. The distinction between the modern derivation and
the historical classification therefore remains essential.

Geach objects that McCall's seven readings do not exhaust the possible
combinations. His additional mixed interpretation translates
\(LAab\) as \(\forall x(Ax\to\Box Bx)\) but \(LIab\) as
\(\exists x\Box(Ax\land Bx)\). Johnson agrees that this is an additional
choice, but notes its failure on Darii LXL. Refuting seven specified
translations is not a proof that no translation can work.

\section{Geach's complex-term system}
\label{sec:p13-johnson-geach}

Geach's G-system uses terms \(\lambda a\) and \(\mu a\), read as
necessarily \(a\) and possibly \(a\). Johnson reports the translations
\(LAab\mapsto Aa\lambda b\), \(LEab\mapsto Ea\mu b\),
\(LIab\mapsto I\lambda a b\), and \(LOab\mapsto Oa\mu b\).
Every term interpretation obeys
\(\lambda a\subseteq a\subseteq\mu a\).
A decision method adds those inclusions as antecedents and applies the
assertoric procedure. This treats modal constructions as complex terms,
rather than applying unrestricted modal operators to every proposition
\parencite[pp.~262--264]{Johnson2004ModalSyllogisms}.

Johnson's figures show why inclusion constraints alone are insufficient.
Darii LXL has \(b\) inside \(\lambda c\) and \(a\) overlapping \(b\),
but the necessary part of \(a\) can miss \(c\).
Cesare LLL can have \(c\) disjoint from \(\mu b\) and \(a\) inside
\(\lambda b\), while \(a\) still overlaps \(\mu c\).
These are failures of a reconstruction, not invalidations of the plain
assertoric Darii and Cesare.

\begin{hbtheorem}{thm:p13-johnson-geach-countercases}{proved=annex:p13.johnson.geach}{Finite complex-term countercases}
Under the stated G-system translations and inclusion constraints,
Darii LXL and Cesare LLL have finite countermodels with all simple and
complex term extensions nonempty.
\end{hbtheorem}

Johnson's complete table~4 compares fourteen moods across three
modal combinations. V means Aristotle accepts the combination on his
reconstruction; Gc means the G-system disagrees. A blank means neither
mark, not that the combination is unassessed.
\parencite[p.~264]{Johnson2004ModalSyllogisms}

\begin{tabular}{@{}lccc@{}}
\toprule Mood & LLL & LXL & XLL\\ \midrule
Barbara & V & V &\\
Celarent & V & V &\\
Darii & V & V,Gc &\\
Ferio & V & V &\\
Cesare & V,Gc & V,Gc &\\
Camestres & V,Gc & V,Gc &\\
Festino & V,Gc & V,Gc &\\
Baroco & V,Gc & &\\
Darapti & V & V,Gc & V,Gc\\
Felapton & V & V &\\
Disamis & V & & V,Gc\\
Datisi & V & V,Gc &\\
Bocardo & V & Gc &\\
Ferison & V & V &\\
\bottomrule
\end{tabular}

The comparison has thirteen disagreements among forty-two cells:
\(29/42\), approximately \(69\%\), agree. The denominator matters:
this is not every modal syllogism, every premise length, or empirical
reasoning performance. Johnson closes this comparison with Kneale's
de-re/de-dicto dilemma and McCall's third approach designed to preserve
the finer distinctions. The mechanisms and limits of that third
approach require their own treatment; they do not follow merely from
the inadequacy of these first two approaches.
